% Source: https://artofproblemsolving.com/community/c3h2164907_algebra_problems
% Stable identity: aops-topic-2164907-q02
Let the circumference be $C$ yards. Moving towards each other, they first meet when together they have covered half a lap, $\dfrac {C}{2}$. Between the first and the second meeting they together cover a full lap $C$, so by the second meeting they have together covered $\dfrac{3C}{2}$, three times as much as at the first meeting. Since the speeds are constant, each of them has covered three times his distance at the first meeting; in particular $B$ has travelled $300$ yards.

At the second meeting $A$ has travelled $C-60$ yards. So
\[300+(C-60)=\dfrac{3C}{2},\qquad\text{hence}\qquad \dfrac {C}{2}=240,\qquad C=\boxed{480}\text{ yards}.\]
(Check: $A$ and $B$ cover $140$ and $100$ yards before the first meeting and $420$ and $300$ yards before the second; $420=480-60$.)
